%% La leyenda de niveles. Va ANTES del índice: PLAN-LIBRO.md §2.5.
\section*{2 · Los dos ejes: quién afirma, y sobre qué}

Con esas palabras ya se puede plantear la pregunta que vuelve en cada línea:
\emph{¿esto es lenguaje objeto o lenguaje meta?} En un libro de metamatemática se
plantea en cada línea, y responderla mal es caro: la inconsistencia que cuenta la
Parte~IV fue exactamente un error de categoría. Por eso aquí el nivel \textbf{se imprime}.

No es una distinción binaria. Hay dos ejes independientes, y cada fragmento de código y cada
enunciado llevan una chapa con los dos:

\noindent\textbf{Eje 1 — ¿quién afirma?} Son cuatro relaciones de derivabilidad distintas, y
\emph{de tipos distintos}: lo que las separa es si arrastran contexto y si son recursivamente
enumerables.

\begin{center}\footnotesize\setlength{\tabcolsep}{4pt}
\begin{tabularx}{\linewidth}{@{}l >{\ttfamily\scriptsize}l cc >{\raggedright\arraybackslash}X@{}}
\toprule
chapa & tipo en Lean & ctx. & r.e. & papel \\
\midrule
\colorbox{nvLean}{\textcolor{white}{\sffamily\scriptsize\bfseries\strut\,Lean\,}} &
  \normalfont--- & --- & --- & la metateoría: un teorema \emph{sobre} el sistema \\
\colorbox{nvDer}{\textcolor{white}{\sffamily\scriptsize\bfseries\strut\,$\Gamma\vdash$\,}} &
  List Formula → Formula → Prop & sí & \textbf{no} & el cálculo $\omega$; sólido en $\mathbb{N}$ \\
\colorbox{nvPrf0}{\textcolor{white}{\sffamily\scriptsize\bfseries\strut\,Prf$_0$\,}} &
  Formula → Prop & no & sí & Hilbert \textbf{intuicionista} \\
\colorbox{nvPrf}{\textcolor{white}{\sffamily\scriptsize\bfseries\strut\,Prf\,}} &
  Formula → Prop & no & sí & Hilbert clásico: \textbf{el que se aritmetiza} \\
\colorbox{nvPrfH}{\textcolor{white}{\sffamily\scriptsize\bfseries\strut\,$\Gamma\,$PrfH\,}} &
  List Formula → Formula → Prop & sí & sí & variante con contexto; ahí vive la deducción \\
\bottomrule
\end{tabularx}
\end{center}

\noindent\textbf{Eje 2 — ¿sobre qué?} Cinco valores.

\begin{center}\small
\begin{tabularx}{\linewidth}{@{}l >{\raggedright\arraybackslash}X@{}}
\toprule
\colorbox{nvLean!18}{\textcolor{nvLean}{\sffamily\scriptsize\bfseries\strut\,Lean\,}} &
  objetos nativos: \ident{Nat}, listas, recursión estructural \\
\colorbox{nvLean!18}{\textcolor{nvLean}{\sffamily\scriptsize\bfseries\strut\,sintaxis\,}} &
  la representación en Lean de la sintaxis: \ident{Term}, \ident{Formula} \\
\colorbox{nvObj!18}{\textcolor{nvObj}{\sffamily\scriptsize\bfseries\strut\,objeto\,}} &
  enunciados de la teoría aritmética: $\obj{+}$, $\obj{\cdot}$, $\obj{\sigma}$, sus axiomas \\
\colorbox{nvCod!18}{\textcolor{nvCod}{\sffamily\scriptsize\bfseries\strut\,código\,}} &
  términos objeto que \emph{denotan} sintaxis: $\cod{\codigo{\varphi}}$, \ident{numeral} \\
\colorbox{nvProv!18}{\textcolor{nvProv}{\sffamily\scriptsize\bfseries\strut\,Prov\,}} &
  dentro del predicado de demostrabilidad: la teoría hablando de sí misma \\
\bottomrule
\end{tabularx}
\end{center}

\noindent \textbf{Los puentes, y el que no existe.} Las cuatro no son intercambiables, y la
dirección importa:
\[
  \Prf_0\,\varphi \;\longrightarrow\; \Prf\,\varphi \;\longrightarrow\;
  \bigl(\forall \Gamma,\ \mathsf{PrfH}\ \Gamma\ \varphi\bigr),
  \qquad
  \Prf\,\varphi \;\longrightarrow\; \text{\ident{axioms}} \vdash \varphi .
\]
El recíproco $\text{\ident{axioms}} \vdash \varphi \to \Prf\,\varphi$ \textbf{es falso, y tiene que
serlo}: si valiera, $\vdash$ sería r.e.\ y Tarski cerraría el paso. Ésa es la asimetría que obliga a
tener dos cálculos en vez de uno, y por eso la chapa distingue en cuál se afirma.
Además \ident{prf0_to_derives} no usa \emph{ninguna} meta-regla $\omega$ y \ident{prf_to_derives}
usa \ident{dne} exactamente una vez: la diferencia entre los dos puentes es, medida,
$\{\ident{FOL.MetaRules.dne}\}$.

\medskip\noindent Tres ejemplos que ocupan tres casillas distintas y que conviene comparar:
\ident{two_mul_consN} afirma \emph{en Lean} sobre \ident{Nat};
\ident{prf_formCode_numeral} afirma \emph{en Prf} sobre \textbf{códigos};
\ident{hFN} afirma lo mismo, pero \emph{en} $\vdash$.
Y \ident{pcc_dot_cons} afirma en \ident{Prf} \textbf{dentro de} \ident{Prov} — cuatro niveles de
distancia entre el primero y el último.

\begin{leccion}[Por qué la chapa es un control, no un adorno]
La cuarta fila del segundo eje es la que se cobró 31 módulos. \ident{tcFn} —«el código del
código»— es una operación sobre \textbf{sintaxis} que se declaró como función \textbf{objeto}, y
una función objeto sólo puede depender de \textbf{valores}. Con la chapa puesta, la mezcla se ve
antes de escribirla.
\end{leccion}
